\documentclass[10pt,a4paper]{amsart}
\usepackage[margin=19mm]{geometry}
\usepackage{amsmath,amssymb,mathtools}
\usepackage[scr=rsfs,cal=euler]{mathalfa}
\usepackage{xfrac}
\usepackage[unicode,hidelinks,bookmarks=false]{hyperref}
\newcommand{\ie}{i.e.\ }
\newcommand{\cf}{cf.\ }
\newcommand{\resp}{respectively\ }
\newcommand{\Subsection}{\subsection}
\providecommand{\nobreakdash}{\nobreak\hbox{-}}
\setlength{\emergencystretch}{2em}
\multlinegap0pt
\usepackage{mathtools}
\usepackage{amssymb}
\usepackage{tikz}
\usepackage{float}
\usepackage{enumerate}
\usepackage[shortlabels]{enumitem}
\newtheorem{prop}{Proposition}
\newcommand{\I}{\mathscr{I}}
\newcommand{\oo}{{\omega}^\omega}
\renewcommand{\subset}{\subseteq}
\title{Cardinal invariants of idealized Miller null sets}
\author{Aleksander Cieślak, Takehiko Gappo, Arturo Martínez-Celis, Takashi Yamazoe}
\date{Source: arXiv:2601.07428v1 (CC BY 4.0)}
\begin{document}
\maketitle

\noindent\emph{Source and licence.} Cardinal invariants of idealized Miller null sets. Version arXiv:2601.07428v1 (CC BY 4.0), distributed by arXiv under the Creative Commons Attribution 4.0 International licence, http://creativecommons.org/licenses/by/4.0/, as recorded in the arXiv licence metadata for this exact version and shown on the abstract page https://arxiv.org/abs/2601.07428v1. Full licence notice: \texttt{licenses/arxiv-2601.07428-CC-BY-4.0.txt}.\par\smallskip

\noindent\textit{Editorial scope.} Counted unit: the article's prop prop:no\_KB\_minimal (source0.tex lines 4156--4158). The article's complete original proof of it is delivered with it (source0.tex lines 4159--4204, one \texttt{proof} environment of 46 lines). No mathematics is written, repaired or re-derived: the statement and the proof are the article's own lines, in the article's own notation, and no retained line is substituted or re-flowed.  No further source-local context is needed: every cross-reference of every retained line resolves inside the two delivered spans.  Nothing was omitted from any retained span, so the package carries no omission record.  No retained line contains a citation, so the wrapper carries no bibliography and no bibliography entry.  No retained line contains a figure environment, an image-inclusion command or drawing code, so no image is delivered and the package declares no asset dependency.  Editorial formatting only, in addition: an AMS \texttt{amsart} preamble that restates the article's own declarations and macros used by the retained lines, namely the environments \texttt{prop} on separate counters, together with the standard packages those lines need.  The retained environments print in this document as Proposition 1 in order of appearance, whereas the article prints the same items with the numbering of its own full document.  The complete native source of the article as distributed in the arXiv e-print for this exact version is bundled unchanged as \texttt{sources/192500/source0.tex}.
    \begin{prop}\label{prop:no_KB_minimal}
	   There is no $\leq_{\mathrm{KB}}$-minimal or maximal member among the family $\{\mathscr{I}_{\mathrm{L}}^f:1\ll f\ll\exp\}$.
    \end{prop}

    \begin{proof}
    	In the rest of the proof, we write $P_n = P_n^{\exp}$. 
    	Let $F\coloneqq\{f\in\oo:1\ll f\ll\exp\}$.
    	For $f,g\in F$ and $n<\omega$, define the following condition:
    	\begin{equation}
    		\label{eq:f_g_pigeon}
    		g(n)^2\leq f(n)\text{ and }g(n)^2 \cdot \left( 1+\sum_{k<n} \left\lceil \frac{2^k}{f(k)} \right\rceil \right) \leq 2^n.
    	\end{equation}
    	It can be easily seen that ``$\forall g\in F~\exists f\in F~\forall^\infty n<\omega~(\text{\eqref{eq:f_g_pigeon} holds})$'' and ``$\forall f\in F~\exists g\in F~\forall^\infty n<\omega~(\text{\eqref{eq:f_g_pigeon} holds})$''.
    	Thus, it suffices to show that if $f,g\in F$ satisfy \eqref{eq:f_g_pigeon} for almost all $n<\omega$, then $\mathscr{I}_{\mathrm{L}}^g\lneq_{\mathrm{KB}}\mathscr{I}_{\mathrm{L}}^f$.
    	Take such $f,g$ and we may assume that \eqref{eq:f_g_pigeon} is true for all $n<\omega$. $\I_L^g \leq_{\mathrm{KB}}\I_L^f$ is obvious as $g\leq^* f$.	
    	To see $\I_L^f \not\leq_{\mathrm{KB}} \I_L^g$, suppose toward a contradiction that $\pi\colon\omega\to\omega$ witnesses $\I_L^f\leq_{KB}\I_L^g$. It suffices to find $A\in\I_L^f$ such that $\pi^{-1}[A]\notin\I_L^g$.
    	To construct such a set $A$, we inductively define $n_i<\omega$ and finite sets $A_i\subset\omega$ as follows.
    	Let $i<\omega$ and assume $(n_j)_{j<i}$ and $(A_j)_{j<i}$ are defined.
    	First, pick $n_i<\omega$ so that if $A_j\cap P_n\neq\emptyset$ for some $j<i$ and $n<\omega$, then for all $m\geq n_i$,
    	\begin{equation}\label{eqn:disjointness}
    		\pi[P_m]\cap P_n = \emptyset.
    	\end{equation}
    	This is possible because $\pi$ is finite-to-one. 
    	To define $A_i$, 
    	put \[N_i \coloneqq \sum_{k<n_i} \left\lceil \frac{2^k}{f(k)} \right\rceil ,\]
    	and let $\langle I_k\rangle_{k< N_i}$ be a subpartition of $\langle P_n\rangle_{n< n_i}$ such that if $I_k \subseteq P_n$ then $|I_k|\leq f(n)$. Put $I_{N_i} = \bigcup_{m\geq n_i}P_m$. The pigeonhole principle implies that there is $k\leq N_i$ such that
    	\[
    	\lvert\pi^{-1}[I_k]\cap P_{n_i}\rvert \geq g(n_i)^2,
    	\]
    	since otherwise \eqref{eq:f_g_pigeon} implies
    	\[|P_{n_i}|=2^{n_i}<g(n_i)^2\cdot (N_i+1)\leq 2^{n_i},\]
    	a contradiction.
    	Let $k_i$ be the least such $k$. Let $\{a^i_j: j<l\}$ be the increasing enumeration of $I_{k_i}\cap\pi[P_{n_i}]$ and let $A_i = \{a^i_j:j<\min\{g(n_i)^2, l\}\}$.
    	
    	Now set $A = \bigcup_{i<\omega}A_i$. Then the following hold:
    	\begin{enumerate}
    		\item For all $i<\omega$, $\lvert\pi^{-1}[A_i]\cap P_{n_i}\rvert \geq g(n_i)^2$.
    		\item For all $i<\omega$, $\lvert A_i\rvert\leq g(n_i)^2$ and if $A_i\cap P_n\neq\emptyset$ for some $n < n_i$ then $\lvert A_i\rvert\leq \lvert I_{k_i}\rvert\leq f(n)$.
    		\item For all $n<\omega$, there is $i<\omega$ such that $A\cap P_n = A_i \cap P_n$ and such $i$ is unique if $A\cap P_n \neq\emptyset$.
    	\end{enumerate}
    	Note that the third clause follows from (\ref{eqn:disjointness}). Then $\pi^{-1}[A]\notin\I_L^g$ because by (1),
    	\[
    	\lvert\pi^{-1}[A]\cap P_{n_i}\rvert \geq \lvert\pi^{-1}[A_i]\cap P_{n_i}\rvert\geq g(n_i)^2
    	\]
    	for all $i<\omega$. To see $A\in\I_L^f$, it suffices to show that $\lvert A\cap P_n\rvert\leq f(n)$ for all $n>0$. By (3), for each $n>0$, there is $i<\omega$ such that $A\cap P_n = A_i\cap P_n$. Then by (2), if $n < n_i$, then $\lvert A_i\rvert\leq f(n)$, or else
    	\[
    	\lvert A\cap P_n\rvert = \lvert A_i \cap P_n\rvert \leq g(n_i)^2 \leq f(n_i) \leq f(n).
    	\]
    	This completes the proof.
    \end{proof}

\end{document}
